\documentclass[12pt,a4paper]{article}
\usepackage[utf8]{inputenc}
\usepackage[spanish]{babel}
\usepackage{amsmath}
\usepackage{amssymb}
\usepackage{listings}
\usepackage{enumitem}
\usepackage{geometry}
\geometry{margin=2.5cm}
\usepackage[utf8]{inputenc} % allow utf-8 input

\usepackage[spanish]{babel}
\usepackage[utf8]{inputenc}

\usepackage[T1]{fontenc}  % use 8-bit T1 fonts
\usepackage{hyperref}    % hyperlinks
\usepackage{url}      % simple URL typesetting
\usepackage{booktabs}    % professional-quality tables
\usepackage{amsfonts}    % blackboard math symbols
\usepackage{nicefrac}    % compact symbols for 1/2, etc.
\usepackage[expansion=false]{microtype}   % microtypography
\usepackage{xcolor}     % colors
\usepackage{graphicx}
\usepackage{subcaption} % Para subfiguras
\usepackage{float}
\usepackage[backend=biber,sorting=ynt,style=apa]{biblatex}


\setlength{\parindent}{0pt} % Eliminar sangría de primera línea en todos los párrafos
% add space between paragraphs
\setlength{\parskip}{1em}

% agregar carpeta imagenes
\usepackage{graphicx}
\graphicspath{ {../imagenes/} }

\addbibresource{../bibliography.bib}

% Configuración para código
\lstset{
 language=Python,
 basicstyle=\ttfamily\small,
 keywordstyle=\color{blue},
 commentstyle=\color{green},
 stringstyle=\color{red},
 numberstyle=\tiny,
 breaklines=true,
 frame=single
}

\title{Aprendizaje Reforzado: Trabajo final}
\author{José Saint Germain}

\begin{document}

\maketitle

\section{Introducción}

El objetivo de este trabajo es realizar generar un análisis y discusión sobre el artículo \textit{A 
reinforcement learning approach to explore the role of social expectations in altruistic behavior} 
\citeyear{castanon2023reinforcement}. En el mismo se realizará un resumen del artículo citado, 
ampliando brevemente sobre los algoritmos en los que se basa, así como en la teoría con la que busca 
discutir.
Adicionalmente, se replicará su trabajo en lenguaje Python para validar sus resultados. Finalmente, se 
generará una discusión sobre las conclusiones a las que arriba, así como posibles extensiones del 
trabajo.

\section{Resumen del artículo}

\subsection{Introducción y Motivación}

El objetivo general del artículo es lograr modelar el comportamiento altruista de individuos. En 
particular, busca llevar al campo computacional la teoría de normas sociales propuesta por Cristina 
Bicchieri (\citeyear{bichieri2006}). Anteriormente, las teorías económicas y psicológicas tenían un 
mayor foco en comportamientos puramente egoístas, es decir, comportamientos que buscan maximizar la 
utilidad individual. Según esta autora, las normas sociales incluyen:

\begin{enumerate}
\item expectativas empíricas: las creencias sobre lo que los demás hacen
\item expectativas normativas: las creencias sobre lo que los demás piensan que se debe hacer
\item preferencias condicionales: la disposición a actuar de acuerdo a sus expectativas
\end{enumerate}

Para llevar a cabo este modelado, los autores recuren a un marco de teoría de juegos, los cuales han 
sido utilizados largamente para capturar comportamiento colaborativo y comportamiento altruista a 
partir de la formación de expectativas de los jugadores. El \textit{Juego del Dictador} resulta 
adecuado para modelar este tipo de comportamiento. Sus reglas básicas son las siguientes:

\begin{itemize}
    \item Dos jugadores: un \textbf{Dictador} y un \textbf{Receptor}.
    \item El Dictador recibe una dotación fija y debe decidir qué fracción donar al Receptor.
    \item El Receptor no tiene poder de decisión; simplemente recibe lo que el Dictador decide donar.
\end{itemize}

Engel (\citeyear{engel2011dictator}) mediante un meta-análisis realizó una revisión literaria de juegos 
del dictador y encontró que, en promedio, los dictadores donan alrededor del 30\% de su dotación, 
entrando en abierta contradicción con la predicción del equilibrio de Nash del juego clásico, que 
indica que los dictadores deberían donar 0\%. Según Bechler y otros (\citeyear{BECHLER2015149}), este 
comportamiento surge de la influencia que tiene la distancia social entre los miembros de este juego.

Finalmente, para encontrar una manera de actualizar las expectativas de los agentes en un juego de
múltiples rondas los autores recurren al aprendizaje por refuerzo, dada la robustez que provee para 
predecir comportamiento prosocial en grandes poblaciones. Este artículo replicará en buena medida
el trabajo realizado por Pereda y otros (\citeyear{pereda2017emergence}) en el cual presenta un modelo 
basado en agentes gobernados por las expectativas generadas por el Juego del Dictador, en las cuales 
los agentes reaccionan a estímulos empíricos basados en la dinámica de aprendizaje reforzado de 
Bush-Mosteller (\citeyear{bush1951mathematical}). La diferencia radica en resolver las siguientes limitaciones encontradas al artículo de Pereda y otros:

\begin{itemize}
\item La ausencia de expectativas normativas.
\item La falta de variabilidad y de dinámicas en la caracterización de los agentes (la susceptibilidad
a estímulos, lo cual puede ser interpretado como la tasa de aprendizajes).
\item El mismo peso asignado a interacciones positivas y negativas.
\end{itemize}

El artículo, entonces, buscará mejorar el modelo para lograr representar los siguientes comportamientos 
humanos observados en el Juego del Dictador: (i) donaciones promedio del 30\%, (ii) la alta 
heterogeneidad en el perfil de los donantes y (iii) la existencia de un número no trivial de humanos 
con un comportamiento no colaborativo.

\subsection{Juego del Dictador}

Dadas las condiciones iniciales y las diferencias establecidas con el trabajo de Pereda y otros, el
Juego del Dictador modelado consta de los siguientes pasos:

\subsubsection{Asignación de aspiraciones}

Todos los agentes (aun no están divididos en dictadores y receptores) comienzan con los siguientes
atributos iniciales:

\textbf{1. Aspiración inicial ($a_{i,t}=0$)}: ubicada aleatoriamente en el rango [0,$\phi$], donde $\phi$ es 
la dotación inicial fija que reciben los dictadores en cada ronda del juego. Este atributo refleja
la cantidad de dinero que esperan obtener en la primera ronda del juego.

\textbf{2. Susceptibilidad ($l_i$)}: representa la sensibilidad del agente a los estímulos sociales. Se 
asigna aleatoriamente en el rango [0,1]. Teóricamente los autores distinguen entre susceptibilidad
normativa y empírica, pero en esta versión del modelo se utiliza una única susceptibilidad.

\subsubsection{Asignación de roles}

Los agentes son catalogados durante esta ronda como dictador (j) o receptor (k) de manera aleatoria. La
cantidad de dictadores y recipientes deben ser la misma para que puedan ser apareados.

\subsubsection{Interacción normativa entre dictadores}

Cada dictador comunica honestamente el porcentaje de la dotación $\theta_{j,t}$ que planea donar al receptor 
a una cantidad $Q$ de otros dictadores. El parámetro $\theta_{j,t}$ no solo toma en cuenta las
aspiraciones originales del dictador sino también un parámetro ruido $\delta_{j,t}^nor$ que toma en
cuenta el efecto de una "mano temblorosa" a la hora de la comunicación. La donación comunicada tiene
la siguiente forma:

$$ \theta_{j,t} = max(0,min(\phi,a_{j,t} * (1 + \delta_{j,t}^{nor}))) $$

La interacción normativa se conforma por la diferencia entre las aspiraciones del dictador y la donación
comunicada a otros dictadores, normalizado por la dotación inicial para obtener un valor entre 0 y 1.

$$ I_{j,t}^{nor} = \frac{\theta_{j,t} - a_{j,t}}{\phi} $$

Con $Q$ cantidad de interacciones, la totalidad de la interacción normativa recibida por el dictador $j$ en la ronda $t$ es:

$$ s_{j,t}^{nor} = \frac{\sum_{m=1}^{Q} I_{m,t}^{nor}}{Q}  = \frac{\sum_{j'}\theta_{j',t}-a_{j,t}\sum_{j'}1}{\phi*Q} = \frac{\sum_{j'}\phi_{j',t}}{\phi*Q}-\frac{a_{j,t}}{\phi}$$

De esta forma, se puede destacar varias cuestiones. Primero, que solo el dictador
que busca consejo actualiza sus aspiraciones en base a la interacción normativa. Segundo,
que el dictador consultado no actualiza las suyas. Tercer, que la interacción $s_{j,t}^{nor}$ 
puede tomar valores entre -1 y 1, siendo positivo cuando la donación comunicada es mayor a la 
aspiración del dictador, y negativo en el caso contrario.

\subsubsection{Los dictadores actualizan sus aspiraciones}

Tanto la susceptibilidad como las aspiraciones de los dictadores se actualizan en base a la interacción normativa 
recibida. Primero, la susceptibilidad se actualiza de la siguiente forma:

% write an if else statement in latex
$$ l_{j,t}^{nor} = 
\begin{cases}
l_{j,t-1}^{nor} * (1+W^{nor,pos} * s_{j,t}^{nor}) & \text{si } s_{j,t}^{nor} \geq 0 \\
l_{j,t-1}^{nor} * (1+W^{nor,neg} * s_{j,t}^{nor}) & \text{si } s_{j,t}^{nor} < 0 \\
\end{cases}
$$

siendo $W^{nor,pos}$ y $W^{nor,neg}$ hiperparámetros (entre 0 y 1) representando los pesos 
asignados a las interacciones normativas positivas y negativas respectivamente.

luego, las aspiraciones se actualizan de la siguiente forma:

$$ a_{j-t} =
\begin{cases}
a_{j,t-1} + (\phi - a_{j,t-1}) * l_{j,1}^{nor} \geq 0 \\
a_{j,t-1} + a_{j,t-1} * l_{j,1}^{nor} < 0 \\
\end{cases}
$$

De estas fórmulas podemos interpretar que el estímulo positivo incrementa la susceptibilidad del agente y,
secuencialmente, sus aspiraciones. Por el contrario, un estímulo negativo reduce la susceptibilidad y las
aspiraciones del agente. Ampliando la interpretación, las interacciones positivas implican que la aspiración
promedio de otros dictadores es mayor a la del agente, llevando al dictadora a aumentar sus aspiraciones 
personales. Lo contrario ocurre en el caso de interacciones negativas.

\subsubsection{Los dictadores proceden a donar}

Cada dictador dona a su respectivo receptor una cantidad $\delta_{j,t}$ con la siguiente fórmula, incorporando
una variable ruido $\delta_{j,t}^{emp}$, que representa la "mano temblorosa" al momento de donar:

$$ \delta_{j,t} = max(0,min(\phi,(\phi-a_{j,t}) * (1 + \delta_{j,t}^{emp}))) $$

Esta donación $\delta_{j,t}$ se convierte en la donación $\pi_{k,t}$ recibida por el receptor $k$ en la ronda $t$.

\subsubsection{Interacción empírica entre receptores}

Como resultado de la donación recibida, cada receptor reacciona al estímulo empírico generado por la misma:

$$ s_{k,t}^{emp} = \frac{\pi_{k,t} - a_{k,t}}{\phi} $$

lo cual se traduce a que, si la donación recibida es mayor a la aspiración del receptor, el estímulo empírico
será positivo, y negativo en el caso contrario.

\subsubsection{Los receptores actualizan sus aspiraciones}

Similar a los dictadores, los receptores actualizan su susceptibilidad y sus aspiraciones en base al estímulo
empírico recibido:

$$ l_{k,t}^{emp} = 
\begin{cases}
l_{k,t-1}^{emp} * (1+W^{emp,pos} * s_{k,t}^{emp}) & \text{si } s_{k,t}^{emp} \geq 0 \\
l_{k,t-1}^{emp} * (1+W^{emp,neg} * s_{k,t}^{emp}) & \text{si } s_{k,t}^{emp} < 0 \\
\end{cases}
$$

$$ a_{k-t} =
\begin{cases}
a_{k,t-1} + (\phi - a_{k,t-1}) * l_{k,1}^{emp} \geq 0 \\
a_{k,t-1} + a_{k,t-1} * l_{k,1}^{emp} < 0 \\
\end{cases}
$$

\subsubsection{Repetir hasta finalizar iteraciones}

Se repiten los pasos anteriores por una cantidad $T$ de rondas, donde en cada ronda se reasignan
los roles de dictador y receptor de manera aleatoria, así como se actualizan las aspiraciones y
susceptibilidades de los agentes.

\subsection{Supuestos de la implementación computacional}

Adicional a las fórmulas y dinámicas presentadas en el artículo, los autores realizan ciertos supuestos
para la implementación computacional del modelo:

\begin{itemize}
    \item La dotación inicial fija $\phi$ es igual a 100.
    \item La variable ruido tanto en la comunicación normativa como en la donación empírica sigue una 
    distribución normal con media 0 y desviación estándar 0.1.
    \item La cantidad de interacciones normativas $Q$ es igual a 5.
    \item La cantidad de rondas $T$ es igual a 5000.
\end{itemize}

\subsection{Resultados Principales}

Las simulaciones muestran que el comportamiento altruista es una propiedad emergente que depende críticamente de la configuración de los pesos de influencia social ($W$).

\begin{itemize}
    \item \textbf{Emergencia de la Norma de Equidad:} Cuando los agentes son sensibles a los estímulos normativos positivos ($W^{nor, pos} > 0$), es decir, cuando aumentan sus aspiraciones al ver que otros son generosos, la población tiende a converger hacia un comportamiento de \textbf{donación alta (cercana al 50\%)}. La norma de equidad se "contagia" y se estabiliza.
    \item \textbf{Colapso hacia el Egoísmo:} Si la influencia normativa es débil o si los agentes pesan desproporcionadamente los estímulos negativos (ej. "si otros dan poco, yo bajo mi estándar"), la población converge hacia el equilibrio de Nash del juego clásico: donaciones cercanas a cero.
    \item \textbf{Dinámica de la Susceptibilidad:} Un hallazgo interesante es que la susceptibilidad ($s$) no necesita ser fija. El modelo permite que $s$ evolucione, mostrando que los agentes pueden volverse más o menos conformistas dependiendo de su historia de interacciones. En escenarios donde la norma es fuerte, la susceptibilidad tiende a mantenerse, reforzando la cohesión del grupo.
    \item \textbf{Papel de las Expectativas Empíricas:} Las expectativas de los receptores actúan como un mecanismo de retroalimentación. Si los receptores se "acostumbran" a recibir poco (sus aspiraciones bajan), esto valida indirectamente el comportamiento egoísta de los dictadores cuando los roles se invierten en rondas futuras.
\end{itemize}

\subsection{Discusión y Conclusiones}

El trabajo concluye que el altruismo en el Juego del Dictador puede explicarse satisfactoriamente mediante un mecanismo de aprendizaje social adaptativo, sin necesidad de asumir preferencias altruistas innatas fijas.

\begin{itemize}
    \item \textbf{Validación de la Teoría de Bicchieri:} Los resultados apoyan computacionalmente la teoría de que las normas sociales son equilibrios sostenidos por expectativas mutuas. El comportamiento prosocial surge de la interacción dinámica entre lo que observamos (normativo) y lo que experimentamos (empírico).
    \item \textbf{Plasticidad del Comportamiento:} El modelo destaca la plasticidad de la moralidad humana. El mismo grupo de agentes puede volverse egoísta o altruista dependiendo simplemente de cómo procesan la información social (los parámetros $W$), lo que tiene implicaciones para el diseño de políticas o intervenciones sociales: para fomentar el altruismo, es crucial visibilizar los comportamientos positivos (reforzar el estímulo normativo positivo).
    \item \textbf{Aporte Metodológico:} El artículo demuestra la utilidad del Aprendizaje por Refuerzo tipo Bush-Mosteller para modelar fenómenos sociológicos complejos, ofreciendo un puente entre la psicología del aprendizaje individual y la dinámica de poblaciones.
\end{itemize}




\printbibliography

\end{document}